\subsection{Primal Heuristics and the Role of Early Incumbents}

Primal heuristics are not an auxiliary convenience in MIP; they are one of the mechanisms through which modern exact solvers become practically effective. Berthold quantified how heavily branch-and-bound performance depends on primal heuristics, especially because improved incumbents immediately affect node pruning and search guidance \cite{berthold2013measuring}. The feasibility pump of Fischetti et al.~\cite{fischetti2005feasibility} remains the canonical example of a relaxation-to-integer heuristic: it alternates between continuous optimization and integer projection in order to move a fractional point toward admissibility. Local branching \cite{fischetti2003local} and relaxation-induced neighborhood search \cite{danna2005exploring} adopt a different philosophy, namely that a high-quality incumbent or relaxation point can define a useful local neighborhood in which exact or heuristic search becomes cheaper. These methods are powerful because they exploit structure already exposed by the relaxation, but they are still fundamentally neighborhood-based methods. When a strong feasible solution lies outside the local basin of the current relaxation point, purely local moves may cycle, stall, or spend many iterations repairing the same violation patterns.

The present work inherits one central lesson from this literature: relaxation information is too valuable to discard. However, it also starts from the observation that local neighborhoods are not always sufficient. HP-PDHG keeps the relaxation at the center of the search, as feasibility pump and local branching do, but augments that center with explicit non-local transitions and staged integrality control. In that sense, it is closer in spirit to a solver-oriented primal heuristic than to a purely black-box metaheuristic, even though its search dynamics are implemented outside a branch-and-cut tree.

At the same time, the method should also be situated relative to the broader intelligent population-search literature that frequently appears in \emph{Expert Systems with Applications}. Classical baselines such as differential evolution \cite{storn1997de}, success-history adaptive differential evolution \cite{tanabe2013shade}, and scalable social-learning particle swarm optimization \cite{cheng2015slpso} are attractive because they provide flexible global search and can operate without explicit derivative information. Their weakness on constrained MIP, however, is that relaxation structure and discrete feasibility geometry are usually represented only through penalties, generic repair, or projection after the fact. Our design tries to retain the exploratory advantages of a population-based heuristic while making the LP relaxation itself a first-class search object. This is the main sense in which HP-PDHG is hybrid: it borrows population diversity from intelligent metaheuristics, but anchors the search in primal-dual relaxation dynamics rather than in a purely black-box fitness landscape.

\subsection{First-Order Relaxation Solvers}

The relaxation backbone of our method comes from first-order primal-dual optimization. The PDHG algorithm introduced by Chambolle and Pock \cite{chambolle2011first} is attractive because each iteration consists of sparse matrix-vector operations and proximal projections, making the method scalable on large sparse saddle-point problems. This style of computation has recently become especially relevant in large-scale linear programming, where Applegate et al.~\cite{applegate2021practical} and Lu and Yang \cite{lu2023cupdlp} showed that carefully engineered primal-dual first-order implementations can be competitive on difficult LP instances. For our purposes, the key trade-off is simple: first-order methods buy cheap scalable iterations, but they do not by themselves explain how to cross the final integrality barrier.

What is missing from this literature, at least for MIP purposes, is a convincing mechanism for converting a high-quality relaxation trajectory into a strict mixed-integer incumbent. Standard PDHG is agnostic to integrality. Even if a trajectory becomes nearly integral, a final hard rounding can destroy linear feasibility. Our method therefore should not be viewed as another attempt to prove first-order optimality results for MIP, but rather as a way of surrounding a first-order relaxation solver with search operators specialized for the feasibility transition that PDHG itself does not address.

\subsection{Learning-Guided MIP Search}

Another important nearby literature concerns learning-guided combinatorial optimization. Khalil et al.~\cite{khalil2016learning} demonstrated that branching behavior in MIP can be learned from data, and subsequent work such as Gasse et al.~\cite{gasse2019exact} showed that graph neural networks can encode branch-and-bound states effectively for exact combinatorial optimization. More recently, neural primal-dual approaches have begun to study how data-driven policies might help generate or improve primal solutions directly \cite{niazadeh2022practical}. These approaches are highly relevant because they underscore a broader point: the weakness of classical heuristics is often not the absence of local information, but the lack of guidance about when and how to depart from the local region currently prescribed by the relaxation.

HP-PDHG is complementary to this line of work. It does not learn a policy from historical instances, but its population, barrier-crossing, and repair structure could naturally accept learned components in future versions. For example, a graph-based policy could prioritize which candidates should receive a non-local perturbation, which variables should be projected more aggressively toward integrality, or which repair actions should be attempted first. We therefore treat learning-guided search not as a competing paradigm, but as a promising future augmentation of the present framework.

\subsection{Barrier-Crossing and Quantum-Inspired Search}

Quantum annealing and adiabatic optimization provide one conceptual origin of our barrier-crossing mechanism. The basic annealing picture is captured by the transverse-Ising interpretation of quantum annealing \cite{kadowaki1998quantum,das2008colloquium}. Experimental realizations on D-Wave hardware triggered a large literature about whether multi-qubit tunneling can provide computational benefits on carefully structured instances \cite{johnson2011quantum,denchev2016computational}. At the same time, the debate over what constitutes genuine quantum speedup remains unresolved and highly problem dependent \cite{ronnow2014defining}. We therefore use this literature cautiously. The relevant lesson for us is narrower and more practical: landscapes with thin but severe barriers may reward occasional non-local transitions that are suppressed by purely local thermal moves.

This view aligns with theoretical and algorithmic work beyond annealing hardware. Crosson and Harrow \cite{crosson2016simulated} showed that simulated quantum annealing can outperform classical simulated annealing on barrier-dominated problems, reinforcing the idea that barrier geometry matters. Hen and Spedalieri \cite{hen2016quantum} studied how constrained optimization can be formulated for quantum annealing without simply dropping the constraint structure. Lucas \cite{lucas2014ising} catalogued Ising encodings for many NP-hard problems, clarifying how discrete optimization can be expressed in quantum-native energy landscapes, while Hadfield et al.~\cite{hadfield2019qaoa} generalized QAOA into the quantum alternating operator ansatz so that hard constraints can be preserved more explicitly. On the variational side, Moll et al.~\cite{moll2018variational} surveyed how near-term quantum optimization algorithms balance expressive search and hardware limitations. Together, these works suggest that constraint handling and search geometry are the key issues, regardless of whether the computational substrate is quantum or classical.

Quantum-inspired classical algorithms occupy the space most directly relevant to our paper. Quantum-inspired evolutionary methods have long used amplitude-like encodings, interference metaphors, and non-classical proposal structures to broaden exploration \cite{das2005quantum,das2011quantum}. Heavy-tailed search processes such as L\'{e}vy flights provide an especially natural surrogate for rare long-range moves because they generate mostly short steps together with occasional large jumps \cite{viswanathan2008levy}. Our barrier-crossing operator draws on that intuition, but differs from generic heavy-tailed randomization in two important respects. First, the jump is embedded inside a PDHG-centered population, so local relaxation information is preserved. Second, the acceptance of a non-local proposal is energy aware and barrier sensitive, rather than a blind restart. This is the main conceptual distinction between our method and earlier quantum-inspired metaheuristics.

\subsection{Position of the Present Work}

Taken together, the literatures above define the niche of HP-PDHG. It is not an exact method like branch-and-cut, not a pure local primal heuristic like feasibility pump or local branching, not a bare first-order relaxation solver, and not a hardware-dependent quantum optimizer. Instead, it is a feasibility-oriented hybrid architecture that combines four ingredients whose interaction is largely absent from prior work: population-based first-order relaxation dynamics, barrier-aware non-local search, staged integrality enforcement, and conservative mixed-integer repair. The method is designed for the regime in which the LP relaxation is informative but insufficient, local repair is too brittle, and full exact search is too expensive to rely on as the sole mechanism for finding incumbents.
